\chapter*{Abstract}
\addcontentsline{toc}{chapter}{Abstract} % add it to table of contents
%\vspace{-1.5em}  % adjust this value as needed
Precisely controlled programs of gene expression are essential for cellular function and organismal viability across all forms of life. Gene regulatory networks (\textsc{GRN}s) provide a high-level mathematical representation of transcriptional regulation by describing the interactions between genes and their regulatory elements, such as transcription factors, across diverse cellular states and contexts. Through this lens, \textsc{GRN}s offer biological insights into processes such as morphogenesis, disease progression, stress adaptation, and response to environmental stimuli. High-throughput profiling of gene expression, particularly single-cell RNA sequencing, has fueled the development of data-driven methods for reconstructing \textsc{GRN}s, sometimes in combination with other data modalities or omics layers. However, in the absence of comprehensive experimentally validated networks corresponding to the exact biological system and experimental conditions under study, evaluating these methods remains a major challenge. %Current \textsc{GRN}-guided simulators adopt coarse assumptions regarding TF coregulation, do not readily preserve gene identities, demand substantial user intervention, and yield benchmarks that deviate from those derived from experimental data, limiting their suitability for this task. 

In this thesis, we address key challenges in benchmarking \textsc{GRN} inference methods through two interconnected studies. \textit{Chapter 3} introduces GRouNdGAN, a \textsc{GRN}-guided simulator that achieves state-of-the-art realism in both steady-state and transient-state scRNA-seq simulations, capturing regulatory properties while preserving gene identities without requiring user intervention.  As a causal generative model, GRouNdGAN expresses genes under the control of their regulators and can sample from both observational and interventional distributions, allowing matched case/control in-silico perturbation experiments. \textit{Chapter 4} formalizes the GRouNdGAN framework as a reproducible protocol, providing in-depth procedural guidance, and introduces GRouNdGAN-toolkit, an extension of the model that includes additional features and customizable modules aimed at facilitating its adaptation to various experimental setups, which are explored through a series of new empirical studies. 

Together, these studies advance the development of \textsc{GRN} inference methods by providing synthetic gold-standard ground truths for any biological system of interest, along with procedures and tools for systematic evaluation, thereby closing the gap between in-silico and experimental data benchmarks of \textsc{GRN} inference methods. 

\newpage

\chapter*{Résumé}
\addcontentsline{toc}{chapter}{Résumé} % add it to table of contents

Des programmes d'expression génique précisément contrôlés sont essentiels à la fonction cellulaire et à la viabilité des organismes dans toutes les formes de vie. Les réseaux de régulation génique (RRG) fournissent une représentation mathématique de haut niveau de la régulation transcriptionnelle en décrivant les interactions entre les gènes et leurs éléments régulateurs, tels que les facteurs de transcription, dans divers états et contextes cellulaires. À travers ce prisme, les RRGs offrent des perspectives biologiques sur des processus tels que la morphogenèse, la progression des maladies, l'adaptation au stress et la réponse aux stimuli environnementaux. Le profilage à haut débit de l'expression génique, en particulier le séquençage d'ARN unicellulaire, a stimulé le développement de méthodes basées sur les données pour la reconstruction des RRGs, parfois en combinaison avec d'autres modalités de données ou couches omiques. Cependant, en l'absence de réseaux complets et validés expérimentalement correspondant au système biologique et aux conditions expérimentales exacts étudiés, l'évaluation de ces méthodes reste un défi majeur.

Dans cette thèse, nous abordons les principaux défis liés à l'évaluation comparative des méthodes d'inférence de RRG à travers deux études interconnectées. Le \textit{Chapitre 3} présente GRouNdGAN, un simulateur guidé par RRG qui atteint un réalisme de pointe dans les simulations scRNA-seq à l'état stationnaire et transitoire, en capturant les propriétés régulatrices tout en préservant l'identité des gènes sans nécessiter d'intervention de l'utilisateur. En tant que modèle générateur causal, GRouNdGAN exprime les gènes sous le contrôle de leurs régulateurs et peut échantillonner à partir de distributions observationnelles et interventionnelles, permettant des expériences de perturbation in silico appariées cas/témoins. Le \textit{Chapitre 4} formalise le cadre GRouNdGAN en un protocole reproductible, fournissant des instructions procédurales détaillées, et présente GRouNdGAN-toolkit, une extension du modèle qui inclut des fonctionnalités supplémentaires et des modules personnalisables visant à faciliter son adaptation à diverses configurations expérimentales, explorées à travers une série de nouvelles études empiriques.

Ensemble, ces études contribuent au développement des méthodes d’inférence des RRGs en proposant des vérités de référence synthétiques applicables à tout système biologique d’intérêt, ainsi que des procédures et des outils d’évaluation systématique, permettant ainsi de réduire l’écart entre les références in silico et expérimentales.